Mediante un ajuste lineal de la gráfica de \( K_{m} \) vs. \( \nu \), se determinó la pendiente, obteniendo la constante de Planck con 
un valor de \( h = \qty{4.00}{\eV\s} \), con un error relativo porcentual de \qty{3}{\percent} respecto al valor registrado en la literatura. 
Además, mediante este mismo ajuste, se obtuvo el valor de la función trabajo del material del ánodo, obteniendo \( \phi = \qty{1.56}{\eV} \),
lo que no coincide con los valores registrados para materiales, siendo el cesio (\ce{Cs}) el metal con el valor de trabajo más bajo, 
con \( \qty{1.9}{\eV} \).

Respecto al comportamiento de las curvas características, se evidenció que la intensidad luminosa \( I \) es proporcional a la magnitud de la
fotocorriente \( I_{p} \), ya que a mayor intensidad luminosa se presenta una mayor cantidad de fotones. Por otro lado, se observó que el 
potencial de frenado es independiente de la intensidad luminosa, ya que se mantuvo constante para cada filtro.
